\begin{mainquestion}{euleridentity}{The Strange Relation Between Exponentials and Sin / Cos}

\questionpoints

% Complex exponentials play a central role in circuit analysis, especially for sinusoidal steady-state behavior.

% \begin{center}
% \begin{circuitikz}
%     \draw
%     (0,0) to[vsourcesin, label=$v_s(t)$] (0,3)
%           to[R, l=$R_s$] (3,3)
%           to[C, l=$C$] (3,0)
%           -- (0,0);
% \end{circuitikz}
% \end{center}

% =========================================================
\begin{partquestion}{complexexponent}
{Some surprising things happen when we allow the time constant in an exponential to be a complex number.  If you don't remember complex numbers, each complex number has a real part and an imaginary part.  The imaginary part consists of a real number multiplied by $\sqrt{-1}$.  Everywhere except EE uses `i' to represent $\sqrt{-1}$.  Since EE generally uses `i' to represent current, we use `j,' so $j=\sqrt{-1}$.  So let's assume we have a function, f, such that:
\[
f(t) = e^{j\omega t}.
\]
Now please compute $\dfrac{df(t)}{dt}$. Relate this answer back to the original function f(t), so your answer should give me a function a, such that a(f(t))=df/dt.}

\begin{answer}[1.7in]
\[
\frac{d}{dt} e^{j\omega t} = j\omega e^{j\omega t}=j\omega~ f(t).
\]
\end{answer}

\end{partquestion}

% =========================================================
\begin{partquestion}{constructg}
{Now let's construct a new function, g, such that
\[
g(t) = \cos(\omega t) + j\sin(\omega t).
\]
Compute $\dfrac{dg(t)}{dt}$. Relate this answer back to the original function g(t), so your answer should give me a function a, such that a(g(t))=df/dt.}

\begin{answer}[1.7in]
\[
\frac{d}{dt}g(t)
= -\omega \sin(\omega t) + j\omega \cos(\omega t)
=j\omega\left(\cos(\omega t) + j\sin(\omega t)\right)=j\omega~ g(t).
\]
\end{answer}

\qnewpage
\end{partquestion}

% =========================================================
\begin{partquestion}{relationship}
{Given these results, what does it say about the relationship between f(t) and g(t)? Show that $f(t)$ and $g(t)$ satisfy the same first-order differential equation and have the same initial condition.}

\begin{answer}[2.5in]
If I have two functions that are the same at t=0, and their derivative is the same at t=0, implies that at t= + or - dt the value will still be the same.  Then using induction (or common sense), it means that the functions must be the same for all time.

Both functions satisfy
\[
\frac{d}{dt}x(t) = j\omega x(t),
\]
and
\[
f(0) = e^{0} = 1,
\qquad
g(0) = \cos(0) + j\sin(0) = 1.
\]
By uniqueness of solutions to first-order differential equations, the two functions must be identical for all $t$, so
\[
e^{j\omega t} = \cos(\omega t) + j\sin(\omega t).
\]
\end{answer}

\noindent
\end{partquestion}

% =========================================================
\begin{partquestion}{cosinesource}
{The voltage source in the circuit produces
\[
v_s(t) = V_0 \cos(\omega t).
\]
We can construct this voltage using the sum/difference of complex exponentials.  To do this we need to use the fact that sin(t) is an odd function of t -- its values for -t are the negative of its values for positive t.  Cos(t) is an even function of time -- its outputs for -t are the same as its values for -t. Using this information and the result from the previous part, write $v_s(t)$ using a complex exponential.}

\begin{answer}[1.2in]
I need to cancel out the sin part of the equation.  Since sin is odd, I just need to add together + and - exponentials.  This leaves 2cos, which is why I divide by 2.
\[
v_s(t) =V_0\frac{e^{j\omega t}+e^{-j\omega t}}{2}.
\]
\end{answer}

\end{partquestion}

% =========================================================
\begin{partquestion}{sinesource}
{The voltage source in the circuit produces
\[
v_s(t) = V_0 \sin(\omega t).
\]
Write $v_s(t)$ using a complex exponential.}

\begin{answer}[1.2in]
I need to cancel out the cos, so I subtract the two exponentials.
\[
v_s(t) =V_0\frac{e^{j\omega t}-e^{-j\omega t}}{2j}.
\]
\end{answer}

\end{partquestion}

\end{mainquestion}
